% ================= CHAPTER 1 =================
\chapter{INTRODUCTION}

\section{Background}
The background reflects a dynamic flow of thought, explaining why a phenomenon encountered can trigger the intention to conduct research. According to \textcite{dewi2022}, the disclosure of research problems should be based on actual phenomena. In this regard, the researcher must feel confident that the phenomenon encountered is still an actual and relevant problem with current conditions \parencite{farhati2022}. This can be obtained by seeking information from the latest literature or consulting with scientific experts related to the discipline concerned \parencite{mallios2010}.

From the researcher's side, the disclosure of this part can be based on the following questions:
\begin{itemize}
    \item What is already known, theoretically or factually, about the problem being studied?
    \item Is there a problem there, are there any doubts found in that problem?
    \item what is interesting about the problem being studied?
    \item Is it technically possible for the problem to be studied?
\end{itemize}
In principle, the research background can answer the question of why the problem is being studied.

\section{Problem Statement}
Problem identification is the elaboration of the problem into several specific sub-problems and is usually formulated in interrogative sentences.

\textbf{Example:} \\
Based on the background description that has been presented, the problem statements raised in this research are as follows:
\begin{enumerate}
    \item How is the structure of the nanocomposite system that can show good two-photon absorption (TPA) efficiency so that this process can be produced at lower light intensities?
    \item What is the mechanism of TPA and the characteristics of the process in the nanohybrid system?
    \item What are the TPA characteristics in semiconductor quantum dots (SQD) modeled more complexly than a two-level system (TLS)?
\end{enumerate}

\section{Research Objectives}
The research objective is the answer to the problem statement. By referring to the example of problem identification formulation, the research objectives are as follows:

\textbf{Example:} \\
This research aims to study the TPA process in SQD-MNP nanocomposites. Specifically, these objectives can be detailed as follows:
\begin{enumerate}
    \item Obtain the configuration of the nanocomposite system structure with good TPA efficiency so that the process can be excited at lower light intensities.
    \item Reveal the principles and mechanisms of TPA as well as the characteristics of the process in the nanocomposite system.
    \item Understand the TPA characteristics in SQD for systems more complex than TLS.
\end{enumerate}

\section{Research Significance}
This section is a sharpening of the specification of the contribution of research results to practical benefits, as well as its scientific contribution to the development of science.

\textbf{Example:} \\
The significance of this research is obtaining information regarding the synthesis method of PP-TiO$_2$ photocatalyst material and the effect of this material on the photodegradation process of organic compounds in peat water. Although this research is still within the scope of a laboratory scale, the results of this research have the potential to be implemented in the water treatment industry sector.
